\indent ಅಲ್ಪ ವಿರಾಮ ಚಿಹ್ನೆ {[},{]} ವಿವರಣಾತ್ಮಕ ಚಿಹ್ನೆ {[}:{]} ಪ್ರಶ್ನಾರ್ಥಕ ಚಿಹ್ನೆ {[}?{]} ಆವರಣ ಚಿಹ್ನೆ ( )\par

\indent ಇವು ಲೇಖನ ಚಿಹ್ನೆಗಳು. ಬರವಣಿಗೆಯಲ್ಲಿ ಈ ಚಿಹ್ನೆಗಳು ಎಷ್ಟು ಪ್ರಾಮುಖ್ಯವೆಂಬುದನ್ನು ತೋರಿಸುವ ಒಂದು ದೃಷ್ಟಾಂತ ಹೀಗಿದೆ—\par

\indent “{\kbvlatin Hang} {\kbvlatin him} {\kbvlatin not}, {\kbvlatin spare} {\kbvlatin him}” “ನೇಣು ಹಾಕಬೇಡಿ ಅವನನ್ನು ಉಳಿಸಿ” ಎಂಬರ್ಥ ಕೊಡುವ ವಾಕ್ಯದಲ್ಲಿ ಅಲ್ಪವಿರಾಮ ಚಿಹ್ನೆಯನ್ನು ಎಲ್ಲಿ ಬಳಸಿದೆ ಎಂಬುದನ್ನು ಗಮನಿಸಿ. ಈ ಚಿಹ್ನೆಯು ಕೈ ತಪ್ಪಿನಿಂದ ಸ್ಥಾನ ಪಲ್ಲಟವಾದಲ್ಲಿ ಎಂತಹ ಪ್ರಮಾದವಾಗುವು ದೆಂಬುದನ್ನು ನೋಡೋಣ.\par

\indent “{\kbvlatin Hang} {\kbvlatin him}, {\kbvlatin not} {\kbvlatin spare} {\kbvlatin him}” “ನೇಣು ಹಾಕಿ ಅವನನ್ನು ಉಳಿಸಬೇಡಿ” ಎನ್ನುವ ತೀರ್ಮಾನಕ್ಕೆ ಕಾರಣ ಅಜಾಗರೂಕತೆಯಿಂದ ಬಳಸಿದ ಅಲ್ಪವಿರಾಮ ಚಿಹ್ನೆ. ದೃಷ್ಟಾಂತವೊಂದರಿಂದಲೇ ಬರವಣಿಗೆಯಲ್ಲಿ ಲೇಖನ ಚಿಹ್ನೆಗಳು ಎಷ್ಟು ಮುಖ್ಯ ವೆಂಬುದು ತಿಳಿದುಬರುತ್ತದೆ, ಆದುದರಿಂದ “ಅವುಗಳ ಬಳಕೆಯ ಬಗೆ ಉತ್ತಮ ತಿಳುವಳಿಕೆ ಇರಬೇಕಾದ್ದು” ಅತ್ಯಾವಶ್ಯಕ.\par

\vspace{4pt plus 4pt minus 1pt}\noindent{\kbvsecfont\fontsize{11pt}{14.5pt}\selectfont\color{kbvaccentdark} (೧)\hspace{0.45em}ಪೂರ್ಣ ವಿರಾಮ ಚಿಹ್ನೆ:}\hspace{0.5em}ವಾಕ್ಯದ ಅಂತ್ಯದಲ್ಲಿ ಪೂರ್ಣ ವಿರಾಮ ಚಿಹ್ನೆಯನ್ನು ಇಡಬೇಕು. ಸಾಮಾನ್ಯವಾಗಿ ಕ್ರಿಯಾಪದದ ನಂತರ ಈ ಚಿಹ್ನೆ ಬರುವುದು.\par

\noindent\hspace*{1.2em}\kbvlabel{ಉದಾ}\hspace{0.6em}ಶ್ರೀರಂಗಪಟ್ಟಣವು ಕಾವೇರಿ ನದಿಯ ತೀರದಲ್ಲಿದೆ.\par

\vspace{4pt plus 4pt minus 1pt}\noindent{\kbvsecfont\fontsize{11pt}{14.5pt}\selectfont\color{kbvaccentdark} (೨)\hspace{0.45em}ಅರ್ಧ ವಿರಾಮ ಚಿಹ್ನೆ:}\hspace{0.5em}ಒಂದು ವಾಕ್ಯವು ಮತ್ತೊಂದು ವಾಕ್ಯಕ್ಕೆ ಪೂರಕವಾಗಿರುವುದುಂಟು, ಅಂತಹ ಸಂದರ್ಭದಲ್ಲಿ ಆ ಎರಡು ವಾಕ್ಯಗಳ ನಡುವೆ ಅರ್ಧ ವಿರಾಮ ಚಿಹ್ನೆ ಬಳಸಬೇಕು.\par

\begin{kbvderiv}
\vspace{4pt plus 3pt minus 1pt}
\noindent{\renewcommand{\arraystretch}{1.24}\setlength{\tabcolsep}{5pt}
\begin{tabularx}{\textwidth}{X X}
ಉದಾ: ಅನ್ನ ಹಳಸಿತ್ತು & ಆದುದರಿಂದ ನಾನು ಊಟ ಮಾಡಲಿಲ್ಲ, \\
ನಾನು ಔಷಧಿ ತೆಗೆದುಕೊಂಡೆ & ಆದರೂ ರೋಗ ಗುಣವಾಗಲಿಲ್ಲ. \\
\end{tabularx}}
\vspace{4pt plus 3pt minus 1pt}
\end{kbvderiv}
